\section{系统启示与未来治理}

\subsection{系统启示板}
Hubert 把 AlphaProof 的设计称为“三层可解释架构”：1）Formalizer/Proof Generator/Verifier 组成 deterministic pipeline；2）Trace/log + dashboard 构成 observability layer；3）Governance review + newsletter 则是 human-in-the-loop layer。AlphaProof 的体验告诉我们：复杂 agent 必须把 pipeline、metric、治理明确分层，才能在 math 这样没灰度的领域安全部署。

\begin{table}[H]
\centering
\begin{tabular}{p{3cm}p{4cm}p{5cm}}
\toprule
层级 & 关注点 & 做法 \\
\midrule
Pipeline & correctness + modularity & Formalizer/Proof Generator/Verifier + Runner 构成 steeped pipeline \\
Observability & metrics & Dashboard + chunk gating + trace log \\
Governance & trust & human review + newsletter + failure curriculum \\
\bottomrule
\end{tabular}
\caption{AlphaProof 的系统启示三层}
\end{table}

\begin{knowledgebox}{系统启示精髓}
把 system design 视为 pipeline + observability + governance 三层架构，可以确保最高敏感的数学推理任务在未知场景也维持可解释与可控。
\end{knowledgebox}

\subsection{未来治理方向}
Hubert 提出“differentiated governance”思路：对不同 search track 设立不同 drift threshold，并把 governance decision 写入 newsletter/backlog。比如 beam search 漏洞由 governance reviewer 直接 patch 相关 policy，而 heuristic track 则由 RL agent 自主重试。他强调：治理不是 judge，而是 feedback policy。

\begin{warningbox}{治理不要简单封杀}
当 governance alert 触发，不是直接 disable agent，而是把 human comment 写回 failure curriculum、提高 reward shaping，再恢复 search；否则会让 system 变成 fragile checkbox。
\end{warningbox}

\subsection{本章小结}
系统启示总结出三层架构，未来治理则通过 differential threshold + feedback loop 让数学 agent 继续拓展边界；这些思考也能反哺其他需要高度可解释的 RL 系统。
\newpage

